% sec8-10-tail.tex -- Sections 8-10 of the C3 paper (draft 2026-11-08).
% Sources: requests-20261106/c3-publication-prior-art-v2.md; notes-20261025/c3-sources.md;
% publish/summaries/c3-*.md; OUTLINE.md section 9; PLAN.md rows G-C3.1-G-C3.3.

\section{Relation to the literature}\label{sec:literature}

\paragraph{The same family.}  Three sources treat the PC or Nullstellensatz degree of
$\mathrm{onto\text{-}FPHP}^{n+k}_n$ with $p\mid k$.
\begin{itemize}
\item Riis's criterion~\cite{riis1993independence} (Theorem~8 of \cite{razborov2001php}) gives
  the upper bound $p^{v_p(k)}$ for every $n$ (Proposition~\ref{prop:riis}).  Our results show
  that it is attained: at the Lean tier for $v_p(k)=1$ from $n=2p-3$ (A), and at the written
  tier for every $k$ from $n=2p^{v_p(k)}-1$ (E).
\item Beame and Riis~\cite{beame1998counting} prove Nullstellensatz degree at least
  $2^\ell-1$ (multiplier degree) for surplus $p^\ell$ and $n\ge((p+2)^\ell-p^\ell)/2$, and at
  most $p^\ell-1$ for $n\ge p^\ell$, and list PC lower bounds for this family as open.  Their designs are the input of D.
\item Riis and Sitharam~\cite[Thm.~33]{riis1997nonconstant} state PC degree at least $2^\ell-1$
  for surplus $p^\ell$ once $N\ge(21+3p^2)(2^\ell-2)-p$.  We have read the statement and the
  preceding lemmas (pp.~36--38, same linear encoding) but not checked the proof, which
  delegates to a tree translation (Lemmas~31, 32) that we have not reconstructed.  We draw no
  conclusion about its correctness, and none from the fact that Razborov's later survey asks
  Question~7 without citing it.  For comparison: at $\ell=2$ it gives $3$ for
  $N\ge2(21+3p^2)-p$, against $p+1$ from $n=2p-1$ (B, Lean) and $p^2$ from $n=2p^2-3$ (F,
  written); at $p=2$, $\ell=3$ it gives $7$ for $N\ge196$, against $5$ from $n=7$ (C, Lean)
  and $8$ from $n=15$ (E, written).
\end{itemize}

\paragraph{Razborov's Question~7.}  The mathematical statements here concern the exact
eventual degree and, for surplus $p^\ell$, growth as $\ell$ varies; under E's named inputs
that degree is $p^\ell$.  We make no claim of resolving the historical formulation of
Question~7 or of priority over Theorem~33 of \cite{riis1997nonconstant}.  None of the Lean
results bears on growth.

\paragraph{Other characteristics and other principles.}  Buss, Grigoriev, Impagliazzo and
Pitassi~\cite{bgip2001} prove linear PC degree for Tseitin formulas mod $p$ in
characteristic $q\ne p$ (when the field contains a root of unity of order $p$) and for counting principles
in other characteristics; Alekhnovich and Razborov~\cite{alekhnovich2003nonbinomial} extend
linear bounds to non-binomial cases; Buss, Impagliazzo, Kraj\'\i\v cek, Pudl\'ak, Razborov
and Sgall~\cite{bikprs1996} prove Nullstellensatz lower bounds for counting principles mod
$q$ over $\mathbb Z_m$ with $q\nmid m$.  All concern the regime in which the characteristic
does not divide the relevant modulus, and none treats the onto principle with $p\mid k$.
Impagliazzo, Pudl\'ak and Sgall~\cite{ips1999} give Razborov's bound for the injective
principle, where the onto and functional axioms are absent and the bound does not depend on
$p$.  The applications of general degree methods~\cite{miksa2015generalized} that we examined do
not yield bounds for the linear onto functional system when $p$ divides $k$, and in the
sources we examined we found no suitable low-degree reduction from the cited counting or
Tseitin principles to this system.

\paragraph{The algebraic inputs.}  The proofs of D, E and F use published results from
representation stability and modular representation theory: FI-modules~\cite{churchellenberg},
matrix loci~\cite{rhoades2024}, stable periodicity~\cite{harman2015},
semisimplification~\cite{hone2025}, chessboard complexes~\cite{shareshianwachs}, Moore spaces
and the Dade group~\cite{yalcin2016}, and endosplit resolutions~\cite{miller2024}.  We make no claim about whether they have been
used in proof complexity before.

\paragraph{Search.}  The searches behind this section (web, ECCC, and the citations of the
sources; October and November 2026), with the list of sources opened and of those read only as
abstracts, are published with each of the findings behind A--C, under \emph{prior art}.

\section{How these results were produced, and what was checked}\label{sec:process}

\paragraph{Who did what.}  All results in this paper were found and proved by AI research
agents (Claude models by Anthropic and Codex models by OpenAI) working in an automated
research harness in October and November 2026.  The author built and ran the harness, chose
the problems, directed the runs, read the outputs and decided what to publish.  He is not a
specialist in this area and does not follow every proof in detail himself.  He publishes the
results as they stand, with the evidence for each stated by its tier, and takes
responsibility for publishing them.  The harness records every claim together with its
evidence, its reviews and the files behind it; those records are the basis of the roles stated
below.  Results A--C come from runs that produce Lean proofs; D--G from one exploratory run
(20261104-0759) with two Claude and two Codex agents, followed by adjudication.

\paragraph{The Lean results.}  The formal statements are supported by kernel-checked proof
terms.  Their interpretation as the claims of this paper also depends on the definitions and
the correspondence established in Section~\ref{sec:defs}.  The library is public, with the exact toolchain and a
manifest; each finding records the commit it was checked at.  Its checks rebuild the proof modules with
the recorded toolchain and then check the published declarations' types and axiom
dependencies with \texttt{\#check} and \texttt{\#print axioms}.  Every theorem uses
only \texttt{propext}, \texttt{Classical.choice} and \texttt{Quot.sound}.  On 6 November 2026 a separate model reproduction rebuilt the named proof targets from the
public repository and checked the published declarations; hardware, timings and memory are
recorded in the ``How to verify'' section of each finding.
For A and B, a second agent in the same run supplied Lean re-proofs, reusing between one and
nine of the first agent's lemmas for A and six for each statement of B.  For C, a second agent
independently checked the level-3 certificate with a program written from the Lean
definitions; this was not a second Lean re-proof.

\paragraph{The written results.}  Each proof of D--G was written by one agent and
reconstructed by at least one other, of the other vendor in most cases; a separate model
referee with no contact with the run re-read the primary sources for each of the three
mechanisms; and the adjudicator Astra (a Codex model outside the run) checked each against the
primary sources and wrote the scoped verdicts we report.  Section~\ref{sec:growth-provenance}
and the corresponding paragraphs of Sections~\ref{sec:periodicity}--\ref{sec:bands} give the
roles per result.  Where reviewers co-developed an argument before reviewing it, the run
records say so.  The paper's later prose reconstructions (Section~\ref{sec:lean},
Appendix~\ref{app:tilting}) and added corollaries were checked in a separate model review of
the paper, which led to corrections.  This is a demanding process, but it is model review,
and we do not present it as more.

\paragraph{What was not done.}  No human expert has reviewed any proof in this paper.  The
external theorems named as inputs were checked at the level of their statements and the
interfaces we use, not re-proved; the exception is \cite[Thm.~8.4.5]{hone2025}, whose proof
was reconstructed, with the result described in Sections~\ref{sec:periodicity}
and~\ref{sec:errata}.  D--G are not formalised.  The finite computations are
exact, but are evidence only for the cases computed.  We would welcome corrections, and will
record them with the published findings.

\section{Open problems}\label{sec:open}

\begin{enumerate}
\item \emph{Formalise E.}  A first step is planned for a revision of this paper: a Lean proof that
  periodicity, and hence the exact degree $p^{v_p(k)}$, follows from the existence of a functor
  with the properties used in Section~\ref{sec:periodicity-oti} or Appendix~\ref{app:tilting}, so that the remaining trust is
  isolated in explicit functor properties and the listed external inputs.  The inputs are deep; a formalisation of F in
  the bands $d<2p$, which needs only the connectivity theorem, is closer.
\item \emph{Complete the proof of \cite[Thm.~8.4.5]{hone2025}} for the semisimplification
  functor, without their Lemma~8.2.1.  This would make Theorem~\ref{thm:periodicity}(b) and
  the threshold $n\ge2q-2$ unconditional.
\item \emph{The uniform spanning statement.}  Prove the spanning statements
  (A)$(n+k,n,r)$, $r=2,3,4$, uniformly over $\mathbb F_3$ for $n\ge7$ and $9\mid k$ ((A) fails
  off the stable range, Section~\ref{sec:bands-reduction}).  Alternatively, formalise
  Lemmas~\ref{lem:hole-avoid} and~\ref{lem:pigeon-ind} and import kernel-checked square
  certificates $S(2)$, $S(3)$, $S(4)$ in the formal spanning definitions; this would discharge
  the hypotheses of the Lean theorem behind Corollary~\ref{cor:f3}.
\item \emph{Thresholds.}  The degree is $p^{v_p(k)}$ from $n=2p^{v_p(k)}-1$ (E) and, for
  $v_p(k)\le2$, from $2p^{v_p(k)}-3$ (A, F).  What is it below?  Over $\mathbb F_2$ with $k=4$
  it is $3$ at $n=3,4$ and $4$ from $n=5$, so the threshold of E is not sharp.
\item \emph{The ideal identity beyond $p^2$.}  Investigate the full identity
  $(I\cap J)_d=(IJ)_d+\sum_{t=1}^d\sigma_t\mathcal A_{d-t}$, the form tested by exact
  computation in the run, for $d\ge p^2$.  For $d<p^2$ the proof of
  Section~\ref{sec:bands-high} establishes the stronger form with $\sigma_1$ and $\sigma_p$
  only; at $d=p^2$ the elementary abelian subgroup is no longer Sylow and the transfer
  argument stops.
\item \emph{Designs of degree $p^\ell-1$.}  Beame and Riis leave the Nullstellensatz degree for
  surplus $p^\ell$ between $2^\ell-1$ and $p^\ell-1$ for odd $p$.  Unit designs of degree
  $p^\ell-1$ would close that gap and, through Theorem~\ref{thm:padding}, give the exact PC
  degree $p^\ell$ on the inputs of Section~\ref{sec:growth} alone.
\item \emph{Riis--Sitharam, Theorem~33.}  Check its proof.
\end{enumerate}
